\SourceLeaf{327}\Needspace{10\baselineskip}

\begin{center}
[2.]

[\emph{Handbuch}, S.\,284.]
\end{center}

Die auf Seite 150 [des Handbuchs] erwähnte Entwicklung lautet wie folgt. Der erste Ausdruck wird zu
\[
= \frac{1}{(01)(23)}\{\sin(02-12)\cos(31-03)+\sin(31-03)\cos(02-12)\}.
\]

Es gilt aber
\[
(12)\sin(02-12) = (01)\sin(01-02)
\]
\[
(13)\sin(31-03) = (01)\sin(01-03),
\]
also
\[
\Omega = \frac{1}{(12)(13)(23)}\{(13)\sin(01-02)\cos(31-03)+(12)\sin(01-03)\cos(02-12)\}.
\]

Ferner gilt
\[
(13)\cos(31-03) = (01)\cos(01-03)-(03)
\]
\[
(12)\cos(02-12) = (02)-(01)\cos(01-02).
\]
also
\[
\begin{aligned}
\Omega ={}& \frac{1}{(12)(13)(23)}\\
&\quad\left\{
\begin{aligned}
&(01)\big[\sin(01-02)\cos(01-03)-\sin(01-03)\cos(01-02)\big]\\
+\;&(02)\sin(01-03)-(03)\sin(01-02)
\end{aligned}
\right\}
\end{aligned}
\]
\[
= \frac{1}{(12)(13)(23)}\{(01)\sin(03-02)+(02)\sin(01-03)+(03)\sin(02-01)\}.
\]

\begin{center}
W. Z. B. W.

[3.]

[\emph{Handbuch}, S.\,201.]
\end{center}

Die oben auf S.\,150 [des Handbuchs] stehenden Sätze über das Viereck lassen sich einfacher wie folgt ableiten.

Die komplexen Größen $p^0, p', p'', p'''$ bezeichnen wie dort die vier Punkte. Die konjugierten Größen seien $P^0, P', P'', P'''$.
